\section{Appendix: Generalization Evidence}
本附录补充三组证据：偏心 flat benchmark、多 patch/disconnected 动态 benchmark 以及 controller statistics。其目的不是改变正文主结论，而是说明当前框架并不只对法向单 patch 场景有效，而是已经能够在偏心加载、多个活跃接触片以及显式 continuity/controller 统计层面给出可复现的结果。

\subsection{偏心 flat benchmark}
偏心 flat benchmark 的重点不在于把法向误差进一步做小，而在于证明当前 native-band 主链已经能够输出可比较的 torque 与 application-point 误差，从而把结论从标量法向力推广到 wrench 层面。表~\ref{tab:eccentric_flat} 给出了 analytic eccentric 与 native-band eccentric 两组结果。可以看到，解析模型上的 torque 指标已经较为稳定，而 native-band eccentric 仍然存在明显 torque/application-point 偏差，这一结果恰好说明：当前主瓶颈已经不再是接触建立时刻，而是更一般动态场景下的 traction consistency 与作用点一致性。

\begin{table}[H]
\caption{偏心 flat 基准中的 wrench 准确性}
\label{tab:eccentric_flat}
\resizebox{\linewidth}{!}{%
\begin{tabular}{lllllllll}
\toprule
Model & Scheme & Force error & Torque error & Peak torque error & Application point error & Impulse error & Energy drift & Release timing error \\
\midrule
analytic\_eccentric & event\_aware\_midpoint\_work\_consistent & 0.004126 & 0 & 0 & 0 & 0.007937 & 0.009755 & 0.002807 \\
native\_band\_eccentric & event\_aware\_midpoint\_work\_consistent & 0.134046 & 0.4375 & 0.4375 & 0.05 & 0.052509 & 0.036002 & 0.007301 \\
\bottomrule
\end{tabular}%
}
\end{table}


\subsection{多 patch / disconnected 动态 benchmark}
多 patch benchmark 用于检验 continuity/controller 在多个分离接触片场景中的稳定性。表~\ref{tab:multipatch_dynamic} 表明，在 analytic multipatch 场景下，当前 work-consistent 动力学器已经能够稳定跨越完整接触相位；而在 native-band multipatch 场景下，active component 数、predictor/corrector Jaccard 与 mismatch fraction 均已可量化地导出。这里最重要的不是误差绝对值已经最优，而是项目现在已经拥有了一个能够显式度量 split/merge 动态的实验接口。

\begin{table}[H]
\caption{多 patch 动态接触基准}
\label{tab:multipatch_dynamic}
\resizebox{\linewidth}{!}{%
\begin{tabular}{lllllllllll}
\toprule
Model & Scheme & Force error & Impulse error & Energy drift & Release timing error & Active component error & Mean active components & Max active components & Mean pred./corr. Jaccard & Mean mismatch fraction \\
\midrule
analytic\_multipatch & event\_aware\_midpoint\_work\_consistent & 0.004126 & 0.007937 & 0.009755 & 0.002807 & 0 & 1.714286 & 2 & -- & -- \\
native\_band\_multipatch & event\_aware\_midpoint\_work\_consistent & 0.249938 & 0.153464 & 0.186517 & 0.007301 & 0 & 2 & 2 & 0.717857 & 0.282143 \\
\bottomrule
\end{tabular}%
}
\end{table}


\subsection{Controller 统计量}
表~\ref{tab:controller_statistics} 总结了当前 controller 在偏心 flat 与 multipatch 场景中的主要统计量，包括 predictor/corrector Jaccard、mismatch fraction、work mismatch、各类 trigger 次数，以及 candidate / recompute cells。它们说明 controller 已经不是黑箱逻辑，而是具有明确可报告行为的统一数值机制。对于后续论文打磨，这一组统计量还可以进一步转化为控制器行为图与附录消融表。

\begin{table}[H]
\caption{Controller 统计与通用性证据}
\label{tab:controller_statistics}
\resizebox{\linewidth}{!}{%
\begin{tabular}{lllllllllllll}
\toprule
Case & Mean pred./corr. Jaccard & Mean mismatch fraction & Max mismatch fraction & Mean work mismatch & Work mismatch triggers & Mismatch triggers & Force jump triggers & Measure jump triggers & Mean candidate cells & Mean recompute cells & Mean active components & Max active components \\
\midrule
native\_band\_eccentric & 0.717857 & 0.282143 & 1 & 0.028588 & 3 & 5 & 8 & 7 & 192 & 166 & 1 & 1 \\
native\_band\_multipatch & 0.717857 & 0.282143 & 1 & 0.005402 & 3 & 5 & 8 & 7 & 140 & 140 & 2 & 2 \\
\bottomrule
\end{tabular}%
}
\end{table}


\section{Appendix: Reference Comparisons and Sensitivity}
本节的目标是补强两类审稿人最关心的问题：第一，native-band 主链与 reference 之间的差距究竟有多大；第二，当前结果是否依赖某组 magic parameter。相应地，本附录给出 reference comparison 与 sensitivity suites 两套自动生成表。

\subsection{参考解对比}
表~\ref{tab:reference_comparison} 把当前统一采用 $\Delta t = 0.01$ 的默认工作流与三类 reference 放在同一个框架下比较：analytic flat 对 exact analytic 和 high-resolution same-scheme，analytic sphere 对 RK4 reference 和 high-resolution same-scheme，native-band flat 对 analytic flat exact 和 high-resolution native-band。该表支持一个重要判断：native-band flat 相对 high-resolution native-band 的误差仍明显小于其相对 analytic flat exact 的误差，这表明当前可见误差中既包含离散误差，也包含模型层差异；因此，后续工作既需要继续改善时间/空间离散，也需要更谨慎地区分“数值逼近误差”和“主链与解析模型之间的物理差异”。

\noindent\textit{For all native-band experiments in this paper, native-band now uses cell-centered volumetric sampling and cell-centered footprint support masking.}

\begin{table}[H]
\caption{参考解对比实验}
\label{tab:reference_comparison}
\resizebox{\linewidth}{!}{%
\begin{tabular}{llllllllllllll}
\toprule
Benchmark & Scheme & Reference & dt & Reference dt & State error & Peak force error & Impulse error & Release timing error & Rebound velocity error & Energy drift & Runtime (s) & Mean candidate cells & Mean recompute cells \\
\midrule
Analytic flat & + work consistency & Exact analytic & 0.04 & -- & 0.236853 & 6.591e-04 & 0.008664 & 7.424e-04 & 0.001131 & 0.001132 & 0.00556 & -- & -- \\
Analytic flat & + work consistency & High-res same scheme & 0.04 & 0.01 & 0.001548 & 0.002362 & 0.010156 & 8.073e-04 & 0.001117 & 0.001132 & 0.00556 & -- & -- \\
Analytic sphere & + work consistency & RK4 reference & 0.05 & 5.000e-04 & 4.986e-04 & 0.001686 & 0.001966 & 8.206e-04 & 7.514e-05 & -7.514e-05 & 0.031156 & -- & -- \\
Analytic sphere & + work consistency & High-res same scheme & 0.05 & 0.02 & 6.981e-04 & 0.001303 & 0.001971 & 7.258e-04 & 8.231e-05 & -7.514e-05 & 0.031156 & -- & -- \\
Native-band flat & + consistent traction reconstruction & Analytic flat exact & 0.06 & -- & 0.18338 & 0.198292 & 0.030332 & 0.075034 & 0.006008 & 0.006026 & 5.983681 & 122.727273 & 106.818182 \\
Native-band flat & + consistent traction reconstruction & High-res native-band & 0.06 & 0.02 & 0.003462 & 0.061817 & 0.035291 & 0.002051 & 0.003035 & 0.006026 & 5.983681 & 122.727273 & 106.818182 \\
\bottomrule
\end{tabular}%
}
\end{table}


\subsection{参数敏感性分析}
表~\ref{tab:sensitivity} 现已按统一 $\Delta t = 0.01$ 重新扫描 grid spacing、band half width、$\eta$、work-mismatch tolerance、continuity 开关以及 consistent traction reconstruction 开关。新的结果更具体地表明：continuity 的主要收益仍然体现在 candidate/recompute 数量下降，而不是把主误差“调小”；$z$ 向分辨率则依然是最敏感的离散方向，其中 $0.005$ 的更细采样在当前窄带配置下反而显著放大了 peak force、impulse 与 energy drift；$\eta$ 的变化也会系统影响主误差；相比之下，band half width 在当前默认设置附近主要改变的是计算开销，而不是标量误差本身。与此同时，由于当前默认 controller 仍采用 \texttt{max\_depth=0}，work-mismatch tolerance 在这组 sensitivity 中几乎不产生可见差异。这些现象对论文写作有两个直接意义：一方面，它们足以支撑“结果并非来自单个特调参数”的论点；另一方面，也更明确地指出 native-band 主链后续最值得继续优化的敏感方向是局部窄带离散与参数稳健性，而不只是继续堆叠新修正模块。

\begin{table}[H]
\caption{敏感性分析}
\label{tab:sensitivity}
\resizebox{\linewidth}{!}{%
\begin{tabular}{llrllllllll}
\toprule
Parameter & Value & Baseline & Peak force error & Impulse error & Energy drift & Release timing error & Runtime (s) & Mean candidate cells & Mean recompute cells & Mean substeps \\
\midrule
baseline & default & True & 0.198292 & 0.030332 & 0.006026 & 0.075034 & 6.812135 & 122.727273 & 106.818182 & 1 \\
grid\_spacing\_z & 0.0200 & False & 0.222404 & 0.035463 & -0.016715 & 0.072853 & 3.518825 & 77.272727 & 77.272727 & 1 \\
grid\_spacing\_z & 0.0100 & False & 0.198292 & 0.030332 & 0.006026 & 0.075034 & 6.778946 & 122.727273 & 106.818182 & 1 \\
grid\_spacing\_z & 0.0050 & False & 0.206797 & 0.035773 & 0.012443 & 0.074066 & 11.34986 & 172.727273 & 152.272727 & 1 \\
band\_half\_width & 8.00 & False & 0.17712 & 0.04189 & 0.009759 & 0.070852 & 5.276463 & 106.818182 & 106.818182 & 1 \\
band\_half\_width & 12.00 & False & 0.198292 & 0.030332 & 0.006026 & 0.075034 & 6.927674 & 122.727273 & 106.818182 & 1 \\
band\_half\_width & 16.00 & False & 0.198292 & 0.030332 & 0.006026 & 0.075034 & 6.802415 & 122.727273 & 106.818182 & 1 \\
eta & 6.00 & False & 0.176774 & 0.038426 & 0.003712 & 0.07341 & 6.484379 & 122.727273 & 106.818182 & 1 \\
eta & 8.00 & False & 0.198292 & 0.030332 & 0.006026 & 0.075034 & 6.585999 & 122.727273 & 106.818182 & 1 \\
eta & 10.00 & False & 0.214222 & 0.020924 & 0.003524 & 0.07314 & 6.676618 & 122.727273 & 106.818182 & 1 \\
work\_mismatch\_relative\_tol & 0.020 & False & 0.198292 & 0.030332 & 0.006026 & 0.075034 & 6.890314 & 122.727273 & 106.818182 & 1 \\
work\_mismatch\_relative\_tol & 0.050 & False & 0.198292 & 0.030332 & 0.006026 & 0.075034 & 6.655358 & 122.727273 & 106.818182 & 1 \\
work\_mismatch\_relative\_tol & 0.100 & False & 0.198292 & 0.030332 & 0.006026 & 0.075034 & 6.587981 & 122.727273 & 106.818182 & 1 \\
continuity\_enabled & False & False & 0.198292 & 0.030332 & 0.006026 & 0.075034 & 7.067634 & 286.363636 & 286.363636 & 1 \\
continuity\_enabled & True & False & 0.198292 & 0.030332 & 0.006026 & 0.075034 & 6.9068 & 122.727273 & 106.818182 & 1 \\
consistent\_traction\_reconstruction & False & False & 0.198292 & 0.030332 & 0.006026 & 0.075034 & 7.168637 & 122.727273 & 106.818182 & 1 \\
consistent\_traction\_reconstruction & True & False & 0.198292 & 0.030332 & 0.006026 & 0.075034 & 7.30704 & 122.727273 & 106.818182 & 1 \\
\bottomrule
\end{tabular}%
}
\end{table}


\section{附录小结}
截至当前版本，本文已经在正文主表之外补齐了三类附录证据：其一，偏心 flat 与 multipatch benchmark 将结论从法向力扩展到 wrench 与多接触片动态；其二，controller statistics 使 continuity/controller 机制具备可解释的行为统计；其三，reference comparison 与 sensitivity suites 为 native-band 主链提供了更完整的说服力边界。它们共同说明：当前项目已经不仅能写出一条主结果线，也已经逐步具备高水平论文所要求的附录证据结构。
